\section{Restricciones de integridad}\label{sec:restricciones}

El MERX ya expresa las llaves de cada conjunto de entidades y las cardinalidades de cada interrelación, y el diseño intuitivo las convierte en llaves primarias, llaves foráneas y columnas obligatorias. Esta sección recoge lo que el MERX no puede expresar: los dominios de los atributos y las reglas del negocio. Cada regla indica de dónde sale y cómo se garantiza. Las que son dependencias funcionales pasan al conjunto $F$ de la sección~\ref{sec:valoracion}.

\subsection{Restricciones de dominio}

{\small\filas
\begin{xltabular}{\linewidth}{L{0.55}L{1.45}}
\rowcolor{azul}\cab{Atributo} & \cab{Valores admitidos} \\
\endhead
Tipo & teórica, práctica o laboratorio \\
Complejidad & bajo, medio o alto \\
Decision & aprobado o nueva confección \\
EstadoA & presente o ausente \\
EstadoS & pendiente, aceptada o rechazada \\
Dia & lunes a sábado \\
PTeorica, PPractica, PLaboratorio & número en $[0, 1]$; las tres suman 1 \\
Cobertura & número en $[0, 1]$ \\
Cantidad, Edad, Capacidad & entero positivo \\
HInicio, HFin & HFin posterior a HInicio \\
FInicio, FFin & FFin posterior a FInicio \\
PGenerar, PValidar & verdadero o falso; al menos uno verdadero en cada autorización \\
\end{xltabular}}

Todos se garantizan con restricciones CHECK sobre la columna o sobre la fila.

\subsection{Reglas del negocio}

Las reglas se agrupan por proceso. La columna \emph{Garantía} dice con qué se hace cumplir: llave o unicidad (UNIQUE), restricción de fila (CHECK), disparador, o la aplicación cuando la regla trata de quién puede hacer una operación y no del estado de los datos.

{\footnotesize\filas
\begin{xltabular}{\linewidth}{>{\bfseries}p{11mm}L{1.35}L{0.75}L{0.6}}
\rowcolor{azul}\cab{Regla} & \cab{Enunciado} & \cab{Origen} & \cab{Garantía} \\
\endhead
\multicolumn{4}{l}{\cellcolor{bandaclara}\textbf{Banco de bloques y catálogos}} \\
RN-01 & Una autorización concede al menos una de las dos capacidades: generar o validar. & RI-03 & CHECK \\
RN-02 & Solo el coordinador que creó un bloque cambia su asignatura, su tipo de actividad o su complejidad. & Enunciado: «sobre aquellos bloques descritos por él» & Aplicación \\
RN-03 & No hay dos bloques con la misma descripción, la misma asignatura y el mismo tipo de actividad. & Enunciado: evitar bloques duplicados & UNIQUE. DF f32 \\
RN-04 & Dos franjas del mismo día no se solapan. En consecuencia, el día y la hora de inicio identifican la franja, y también el día y la hora de fin. & Decisión de la sección~\ref{sec:requerimientos} & Disparador; UNIQUE. DF f35 y f36 \\
RN-05 & Dos semestres no se solapan. La fecha de inicio identifica al semestre, y también la de fin. & Decisión & Disparador; UNIQUE. DF f33 y f34 \\
\multicolumn{4}{l}{\cellcolor{bandaclara}\textbf{Generación, validación y revisión}} \\
RN-06 & Quien genera un horario tiene autorización con PGenerar en el programa del horario. & Enunciado: «autorizado a generar» & Disparador \\
RN-07 & Quien valida un horario tiene autorización con PValidar en el programa del horario. La validación no es anterior a la creación del horario. & Enunciado: «autorizado a validar» & Disparador \\
RN-08 & Una sesión solo programa bloques que abordan asignaturas del programa del horario. & Decisión & Disparador \\
RN-09 & En un horario, cada franja tiene a lo sumo una sesión. & Enunciado: impedir la asignación doble & UNIQUE (\#H, \#F). DF f37 \\
RN-10 & En un mismo semestre, dos sesiones de horarios finales distintos no usan la misma aula en la misma franja. & Ocupación de aulas & Disparador al programar la sesión y al aprobar el horario \\
RN-11 & La revisión no es anterior a la creación del horario. Un horario es final si y solo si su revisión tiene decisión aprobado. & RI-10 & CHECK; disparador \\
RN-12 & Un horario sustituye a otro solo si la revisión de este decidió nueva confección. Los dos son distintos y tienen el mismo programa y el mismo semestre. & RI-11 & Disparador \\
\multicolumn{4}{l}{\cellcolor{bandaclara}\textbf{Asignación y asistencia}} \\
RN-13 & Solo se asigna un horario final, y solo a un estudiante matriculado en el programa del horario. & Enunciado: horarios que «le son asignados» & Disparador \\
RN-14 & Un estudiante no tiene asignados dos horarios del mismo semestre con sesiones en la misma franja. & Enunciado: impedir la asignación doble & Disparador al asignar y al programar una sesión \\
RN-15 & Solo registra asistencia un estudiante asignado al horario de la sesión. La fecha cae dentro del semestre del horario y en el día de la semana de la franja. & Enunciado: asistencia «ante un horario impartido» & Disparador \\
\multicolumn{4}{l}{\cellcolor{bandaclara}\textbf{Solicitud de cambio}} \\
RN-16 & Cuando el estudiante presenta la solicitud, el horario de origen está asignado a él. & RI-16 & Disparador \\
RN-17 & Una solicitud pendiente no tiene fecha de resolución ni destino. Una aceptada tiene ambos. Una rechazada tiene fecha de resolución y no tiene destino. La resolución no es anterior a la solicitud. & RI-16 & CHECK \\
RN-18 & El destino es un horario final, distinto del origen, del mismo programa y del mismo semestre. & Decisión & Disparador \\
RN-19 & El coordinador de apoyo es distinto del que atiende la solicitud. & Enunciado: «adicional al que tienen asignado» & CHECK \\
RN-20 & Al aceptarse la solicitud, el estudiante deja el horario de origen y queda asignado al de destino en la misma transacción. Se siguen cumpliendo RN-13 y RN-14. & Enunciado: «el resultado del cambio manual» & Disparador \\
\multicolumn{4}{l}{\cellcolor{bandaclara}\textbf{Consulta e histórico}} \\
RN-21 & Un coordinador solo consulta horarios de programas en que está autorizado. Un estudiante solo consulta los horarios que tiene asignados. & Enunciado & Aplicación \\
RN-22 & Los reportes de bloques más impartidos y de equilibrio usan solo horarios finales. & Demandas 21 y 24 & Consulta \\
RN-23 & El histórico no comparte tablas ni llaves foráneas con el modelo operacional. & RNF-02 & Esquema separado \\
\end{xltabular}}

RN-10 y RN-14 comparan sesiones dentro del semestre del horario. Como la franja es semanal, dos horarios de semestres distintos pueden usar la misma aula en la misma franja, y un estudiante puede tener la misma franja ocupada en semestres distintos. Un horario sustituido o todavía sin aprobar no ocupa aulas para RN-10 y no se puede asignar.

\subsection{Integridad referencial}

Cada llave foránea de la tabla de la sección~\ref{sec:relacional} exige que el valor exista en la relación referenciada. Ninguna admite nulo, salvo \#HDes y \#CApo en Solicitud, que corresponden a los extremos (0,1) de DESTINO y APOYAR. La llave foránea (\#H, \#B) de Asistencia apunta a Sesión, de modo que solo hay asistencia a sesiones programadas. Las relaciones Asistencia, Validación y Revisión no se borran en cascada: guardan la historia del horario y del estudiante.
